\section{Introduction}
Infectious disease dynamics are heavily influenced by the underlying contact network of individuals.
In standard epidemic models, the network structure is often assumed to be static.
However, real-world populations exhibit adaptive behavior where individuals modify their contacts in response to the epidemic.
This project investigates a Susceptible-Infected-Recovered (SIR) model on an adaptive network, where susceptible individuals break links with infected neighbors and form new connections to non-infected agents.
The feedback loop between the disease spread and the evolving network topology makes the likelihood function for this model intractable \parencite{demirel2017dynamics}.
Standard statistical methods, such as Maximum Likelihood Estimation, are therefore difficult to apply directly to such complex stochastic systems.
To address this challenge, we employ Simulation-Based Inference (SBI) to estimate the core parameters: the infection probability $\beta$, the recovery probability $\gamma$, and the rewiring probability $\rho$.
Since SBI requires a significant number of forward simulations, we implemented a high-performance simulator using Python and Numba to ensure computational efficiency during the inference process.
We are provided with observed data from $R=40$ independent realizations of the stochastic process.
The available measurements include time-series data for the infected fraction and rewiring counts, as well as final degree histograms at the conclusion of the simulation.
The goal of this report is to describe the implementation of Approximate Bayesian Computation (ABC) and explore advanced sequential methods to obtain robust posterior distributions for the unknown parameters.